\chapter{Pruebas y validación}
\label{cap:pruebas}
\section{Plan de pruebas}
\pendiente{Vincular requerimientos, casos, entornos y criterios de aceptación.}
\subsection{Casos del incremento rítmico R6}
El plan del incremento descrito en la sección~\ref{sec:implementacion-ritmo-r6}
verifica las acciones de la sesión pública con tiempo, audio y almacenamiento
controlados. Comprende dos resultados altos, omisiones, silencios con golpes
extra, mezcla a tempo inferior, práctica y repetición, techo de tempo,
reapertura, revisiones tardías, guardados concurrentes, exportación y reinicio.
Los fallos y reintentos de escritura se comprueban sin atribuir progreso a
resultados todavía pendientes. Un banco sintético independiente permite
inspeccionar las decisiones sin modificar el perfil real.

La ejecución en Android debe comprobar además los controles, el texto
ampliado, el acceso al panel y el recorrido de dos ejercicios con desbloqueo
y repetición. Estas pruebas se prepararon pero no se ejecutaron en esta sesión
por solicitud del usuario de no abrir el emulador. Las pruebas JVM y las
simulaciones no sustituyen la comprobación del dispositivo ni una evaluación
con estudiantes. Los resultados y límites del host se registran en la guía
del incremento y su evidencia, separadamente de esas verificaciones pendientes.
\subsection{Casos del incremento de Defender D3}
El plan correspondiente a la sección~\ref{sec:implementacion-defender-d3}
comprueba enfriamiento, límites de viaje, conservación de plazos de enemigos
visibles y cómputo de pérdidas e intentos globales. La sesión pública se
ejercita con reloj y audio controlados. Para persistencia se verifican
idempotencia, concurrencia, lectura suspendida, datos incompatibles,
exportación y recuperación junto a los datos de ritmo.

Un caso específico finaliza aparentemente la partida por tercera llegada,
entrega después una respuesta con timestamp previo al plazo y abandona la
ejecución reabierta. El resultado confirmado debe ser una interrupción y no
debe sumar progreso. Los ocho escenarios de inspección son sintéticos y no
se escriben en el perfil real. Las pruebas Android de desbloqueo, reapertura
y simulación quedaron preparadas y compiladas, sin ejecución por solicitud
del usuario de mantener cerrado el emulador. La evidencia seleccionada de host
se conserva en \texttt{docs/evidencias/defender-d3/}; las medidas físicas y la
validación con estudiantes continúan pendientes.
\subsection{Casos de tríadas de Defender D4}
Los casos de la sección~\ref{sec:implementacion-defender-d4} verifican el
catálogo de voces, intervalos, grafías y límites de pentagrama. La sesión
pública comprueba respuesta completa, rechazo de raíz enarmónica distinta,
conservación tras error, limpieza tras acierto, unicidad de objetivos y
timestamp con calidad capturada. La prueba de síntesis examina las tres
frecuencias de cada calidad sobre Do natural y los límites PCM; no constituye
una medición acústica del dispositivo.

Los casos de persistencia y ViewModel verifican progresos separados, dos
resultados altos con cierre confirmado, reapertura, exportación, recuperación,
reinicio y lectura de registros del formato previo. La ruta completa de
calidades se genera con sesiones sintéticas aisladas del perfil. La evidencia
seleccionada se conserva en \path{docs/evidencias/defender-d4/}.
Dos casos Android preparados comprueban el recorrido de tríadas y la
simulación aislada; su ejecución, la legibilidad con cuatro objetivos y
texto ampliado, el sonido físico y TalkBack quedan pendientes por el pedido
de mantener cerrado el emulador. No se infieren resultados de aprendizaje
desde estas pruebas de software.
\subsection{Casos de séptimas e inversiones de Defender D5}
El plan de la sección~\ref{sec:implementacion-defender-d5} recorre las
combinaciones de calidad, inversión y clave, comprobando altura, grafía,
orden de voces, disposición cerrada y límites de líneas adicionales.
También verifica el desplazamiento de cabezas contiguas, la respuesta por
raíz frente a un bajo diferente y la unicidad entre objetivos vivos.
El banco PCM examina las voces de los extremos de cada contexto de clave
simple con señales sintéticas; no mide el altavoz ni la latencia del teléfono.

Una ruta aislada de 168 partidas recorre los 84 objetivos mediante la sesión
pública, con reloj y almacenamiento controlados. Comprueba la alcanzabilidad
del grafo y la evidencia por clave sin modificar el perfil real.
Otros casos verifican que una inversión diferente o una mezcla sin fa no
otorguen progreso, así como lectura de perfiles anteriores, reapertura y
exportación de séptimas invertidas. La evidencia seleccionada se conserva
en \texttt{docs/evidencias/defender-d5/}.
Los recorridos Android se prepararon para séptimas en fa, mezcla, pausa y
reapertura; solo se compilan por el pedido de mantener cerrado el emulador.
La legibilidad con cuatro objetivos, el texto ampliado, la reproducción
física y la evaluación con estudiantes permanecen pendientes.
\section{Verificación del núcleo de intervalos vocales}
La verificación de V1 utiliza la sesión pública: acciones de inicio y cierre,
estimaciones de tono y tiempos controlados producen estados y resultados.
Los casos comprueban relaciones ascendentes y descendentes, unísono,
traslado de registro, conservación de errores de octava y rechazo de
continuidades falsas por silencio o huecos temporales. También se controla
que una interrupción conserve las mediciones sin conceder un resultado musical.

Un recorrido adicional genera PCM sinusoidal y lo procesa mediante el YIN
existente antes de enviarlo a la sesión. Sus resultados corresponden a señales
sintéticas en JVM; no miden voz humana, captura Android ni latencia acústica.
Los reportes del incremento se identifican en
\texttt{docs/evidencias/intervalos-vocales-v1/}. La conexión con la interfaz y
la evaluación sobre dispositivos y estudiantes siguen pendientes para V1.

\subsection{Plan y verificación de la ronda vocal V2}
La integración se verifica por acciones públicas del modelo, con reloj,
reproducción y captura controlados en sus límites. Los casos recorren una
ronda completa y comprueban que la segunda captura espere la liberación de
la primera. Se incluyen cancelaciones repetidas, callbacks anteriores al
cierre, timeout entre notas, fallos de audio y rechazo de una nueva ronda
que descartaría una primera nota ya adquirida.

Dos pruebas Android se preparan para la navegación del libro y para el
resumen con entradas sintéticas. Su compilación no acredita ejecución ni
funcionamiento acústico. Quedan pendientes la revisión a 411, 360 y 320 dp,
texto ampliado al 200\%, permisos y ciclo de vida en dispositivo, TalkBack
y voces reales. La evidencia del incremento se identifica en
\texttt{docs/evidencias/intervalos-vocales-v2/}; no se interpreta como
validación pedagógica.

\subsection{Persistencia y recuperación vocal V3}
Las pruebas atraviesan las acciones públicas del modelo, el historial y las
salidas de exportación. Se verifica que una ronda completa sobreviva a reabrir,
que varias cancelaciones conserven un único resultado y que los reintentos de
una confirmación incierta no dupliquen el identificador. Los casos de datos
dañados y esquema no compatible comprueban la conservación del archivo original
y su inclusión en la copia de recuperación junto al resultado pendiente.

También se verifica el drenaje previo al reinicio, la independencia respecto
del laboratorio de voz y la conservación del pendiente cuando termina una
lectura anterior. Se preparan dos pruebas Android adicionales para reapertura
del historial y fallo/reintento; solo se compilan. La evidencia se identifica en
\texttt{docs/evidencias/intervalos-vocales-v3/}. Estas comprobaciones no evalúan
voz real, latencia acústica ni eficacia pedagógica.

\subsection{Adaptación y progresión vocal V4}
Las pruebas verifican cinco rondas enlazadas, mantenimiento del criterio
durante el sostenido, hitos por intervalo y dirección, orden de contenidos
y conservación de accesos adquiridos. Se comprueban error de octava,
dispersión, protocolo abreviado, series parciales, tolerancia ajena, ordinales
repetidos y confirmaciones de guardado inciertas. Un recorrido sintético de
250 rondas atraviesa las sesiones públicas y el commit hasta registrar hitos
para los 25 objetivos firmados. Este recorrido caracteriza las reglas del
software y no constituye evidencia de aprendizaje.

Dos pruebas Android adicionales se preparan para cancelación/simulación y
una serie de cinco rondas con entradas controladas. Solo se compilan por la
restricción de no emular. La evidencia está identificada en
\texttt{docs/evidencias/intervalos-vocales-v4/}; quedan pendientes legibilidad,
TalkBack, voz real y calibración de parámetros.

\subsection{Ancla dada V5}
La verificación recorre reproducción de una única referencia, captura exclusiva
de la nota objetivo y persistencia de su procedencia. Se comprueban dirección,
unísono, error de octava, silencio, timeout y el pipeline de PCM sintético a YIN
y sesión. Los casos de progreso verifican prerrequisitos por dirección, separación
de tolerancias y rachas, rechazo de referencias incoherentes y compatibilidad
con snapshots V4. Una ruta sintética de 500 rondas cubre los 25 objetivos de
imitación y los 25 de ancla dada, sin participación de estudiantes.

Dos pruebas Android adicionales se preparan para el flujo de una sola captura
y para el bloqueo de acceso automático sin hito de imitación. Solo se compilan.
La evidencia se identifica en \texttt{docs/evidencias/intervalos-vocales-v5/}.
Voz real, legibilidad nativa, TalkBack y calibración siguen pendientes.

\subsection{Producción libre V6}
La verificación comprueba el inicio sin reproducción, las dos capturas separadas,
el ancla elegida por el usuario y la ausencia de referencias sonoras inventadas
en el historial. Se ejercitan prerrequisitos por dirección, tolerancias y rachas
independientes, series mezcladas o abandonadas, timeout entre notas, segundo
plano y callbacks tardíos. La compatibilidad con snapshots V5, el reintento de
confirmación incierta, la reapertura y las salidas JSON y ZIP se verifican sin
duplicar progreso.

Una ruta sintética de 750 rondas atraviesa los 25 objetivos de cada uno de los
tres modos, mediante sesiones públicas y almacenamiento aislado. Es una prueba
de reglas del software, sin estudiantes ni medición acústica. Dos pruebas
Android adicionales se preparan para el flujo silencioso con historial y el
bloqueo sin hito de ancla dada; solo se compilan por la restricción de no emular.
La evidencia se identifica en \texttt{docs/evidencias/intervalos-vocales-v6/}.
Voz real, legibilidad nativa, TalkBack y calibración continúan pendientes.

\section{Resultados obtenidos}
\pendiente{Incorporar resultados comprobables con versión, fecha, datos, dispositivo y condiciones.
Separar pruebas automatizadas, simulación, medición física y evaluación con estudiantes.}
\section{Métricas de calidad}
\pendiente{Definir métricas y métodos antes de interpretarlas. La demostración funcional
por sí sola no acredita precisión acústica ni eficacia pedagógica.}
